\subsection{Блок-схема численного метода}

На вход реализуемого метода подаются следующие данные:

\begin{itemize}
  \item \( f \) --- номер интерполируемой функции; является целым числом.
  \item \( a \) --- левая граница интерполируемого интервала; является вещественным числом двойной точности.
  \item \( b \) --- правая граница интерполируемого интервала; является вещественным числом двойной точности.
  \item \( x \) --- значение аргумента для интерполяционного многочлена; является вещественным числом двойной точности.
\end{itemize}

На выходе метода ожидается значение интерполяционного многочлена в точке \( x \), которое является вещественным числом двойной точности.

Далее приводятся блок-схемы, реализующие численный метод:

\begin{scheme}
  \centering
  \begin{tikzpicture}[gost flowchart, start chain=flow going below]
    \node [terminator, on chain] (start) {Начало};
    \node [terminator, on chain, join] (end) {Конец};
  \end{tikzpicture}
  \caption{Вычисление значения интерполирующей функции \( F(x) \) заданной точности.}
\end{scheme}

\begin{scheme}
  \centering
  \begin{tikzpicture}[gost flowchart, start chain=flow going below]
    \node [terminator, on chain] (start) {Начало};
    \node [terminator, on chain, join] (end) {Конец};
  \end{tikzpicture}
  \caption{Построение интерполяционного многочлена в форме Лагранжа \( L_n(t) \).}
\end{scheme}

\begin{scheme}
  \centering
  \begin{tikzpicture}[gost flowchart, start chain=flow going below]
    \node [terminator, on chain] (start) {Начало};
    \node [terminator, on chain, join] (end) {Конец};
  \end{tikzpicture}
  \caption{Поиск корней многочлена Чебышева первого рода \( T_n(t) \).}
\end{scheme}

\begin{center}
\begin{tikzpicture}[%
    start chain=going below,    % General flow is top-to-bottom
    node distance=6mm and 60mm, % Global setup of box spacing
        ]
        \node [cloud] (start) {Начало};
        \node [block, join] (phase1) {Этап 1};
        \node [cyclebegin, join=by red] (phase2) {Этап 2};
        \node [block, join=by green] (phase3) {Этап 3};
        \node [cycleend, join] (phase4) {Этап 4};
        \node [subroutine, join, subrtshape w = 3mm, fxd] (phase5) {Этап 5 1 2 3 4 5 6 7 8
9 0};
        \node [io, join, fxd] (input) {Этап 6 - ввод данных};
        \node [block, join] (phase7) {Этап 7};
        \node [decision, join] (condition) {Условие};
        \node [connector] (finish) {Конец};
        \node [block, left of = phase4, node distance = 3cm] (correction) {Коррекция};
        \node [print_block, right of = phase4, node distance = 3cm] (print) {Print};
        \path [line, red] (condition) -| node [u,near start] {Нет} (correction);
        \path [line] (correction) |- (phase1);
        \path [line] (phase2) -| node [r,near end] {Печать} (print);
        \path [line, green] (condition) to node [r] {Да}(finish);
\end{tikzpicture}
\end{center}